\section{Two labels: path-ordered products and the maximal normal form}
\label{sec:positive-descendents}

Proposition~\ref{prop:first-decoration} compares the GKT transition with the
diagram $\mathfrak D$ one ray at a time.  This section compares
the two along a path, first for two labels of twists $(1,0)$ and $(0,1)$,
where $\mathfrak D$ has a mixed wall, and then for two labels of the same
twist $(1,0)$, where it has none.  As in
Section~\ref{sec:gkt-decoration}, take two labels with
\[
 t_1=u_{i,1},\quad\tau(i)=(1,0),
 \qquad
 t_2=u_{j,1},\quad\tau(j)=(0,1),
\]
and work over $R_2=\QQ[t_1,t_2]/(t_1^2,t_2^2)$.  Operators act on
polynomials in $x,y$ over $R_2$ through the vector fields
\eqref{eq:GKT-generator}, and the rightmost factor of a product acts first.
Put
\begin{equation}\label{eq:two-label-flows}
 \Sigma_1=\exp\!\left(-\tfrac12t_1X_{1,0}\right),\qquad
 \Sigma_2=\exp\!\left(-\tfrac12t_2X_{0,1}\right),\qquad
 Z_c=\exp\!\left(c\,t_1t_2X_{1,1}\right)\quad(c\in\QQ).
\end{equation}
Under Theorem~\ref{thm:critical-embedding}, $\Sigma_1$ and $\Sigma_2$ correspond to the
incoming singleton factors $e^{\mathsf A}$ and $e^{\mathsf B}$.  Each $Z_c$
is central in the group generated by these operators, and
$Z_cZ_{c'}=Z_{c+c'}$.  Of the two balanced monomials $t_1t_2xy^5$ and
$t_1t_2x^5y$ of the support $J=\{i,j\}$, the extremal GKT potentials
$W^{\min}$ and $W^{\max}$ of \eqref{eq:mixed-potentials} keep only the one of
least, respectively greatest, $x$-degree; its coefficient
$12=(-1)^{|J|-1}\nu$ is determined by the GKT invariant $\nu=-12$ of $J$.
Direct differentiation gives
\begin{equation}\label{eq:two-label-derivatives}
\begin{gathered}
 X_{1,0}(x^4+y^4)=4x^5-8xy^4,\qquad X_{1,0}(y^5)=-10xy^5,\\
 X_{0,1}(x^4+y^4)=8x^4y-4y^5,\qquad X_{0,1}(x^5)=10x^5y,\\
 X_{1,1}(x^4+y^4)=8x^5y-8xy^5.
\end{gathered}
\end{equation}

\subsection{Path-ordered products near the joint}

Let $(s_1,s_2)$ be affine coordinates on the chart of $\Delta$ drawn in
Figure~\ref{fig:cofactor-supports}(a), with the joint $O$ at the origin.  By \eqref{eq:cofactor-singletons} the two singleton
chords are the lines $3s_1+4s_2=0$ and $3s_1+2s_2=0$.  By
\eqref{eq:outgoing-cofactor} with $c_1=c_2$, the first mixed wall is the
segment from $O$ to $B$ on the line $s_1+s_2=0$, $s_2\geq0$.  Modulo $(t_1^2,t_2^2)$ these are the only walls
through $O$ with a nontrivial factor.  With the coorientations of
Section~\ref{sec:jacobi-theta}, the three walls carry the following factors:
\begin{center}
\begin{tabular}{lll}
\toprule
Support & Positive normal & Attached factor\\
\midrule
$3s_1+4s_2=0$ & $(3,4)$ & $\Sigma_1$\\
$3s_1+2s_2=0$ & $(3,2)$ & $\Sigma_2$\\
$s_1+s_2=0$, from $O$ to $B$ & $(1,1)$ & $Z_{-3/4}$\\
\bottomrule
\end{tabular}
\end{center}
A path crossing a wall in the direction of its positive normal applies the
attached factor, and otherwise its inverse
(Definition~\ref{def:Jacobi-wall}).  Choose $\varepsilon>0$ so small that
the disk of radius $2\varepsilon$ about $O$ meets no wall whose factor is
nontrivial modulo $(t_1^2,t_2^2)$ other than these three and contains no
singular point; this is possible by Proposition~\ref{prop:proper-supports}.  The
transitions at the singular points are carried by cuts running from the
singular points outward (Section~\ref{sec:jacobi-theta}), so they do not meet
this disk, and the regions of $\mathfrak D$ adjacent to $O$ share one frame.  Let $\gamma_+$
and $\gamma_-$ run from $\ell=(-\varepsilon,0)$ to $r=(\varepsilon,0)$ along
the horizontal segments $s_2=\varepsilon/10$ and $s_2=-\varepsilon/10$,
respectively, joined to $\ell$ and $r$ by vertical segments, and write
$\Theta_{\gamma_\pm}$ for the path-ordered products.  Along $\gamma_+$ the
walls are met in the order of increasing slope of their GKT generators, so
$\ell$ lies in the region preceding all three walls in that order, and we
place $W^{\min}$ there.

\begin{proposition}[Two-label comparison]
\label{prop:two-label-transport}
\begin{enumerate}
\item The path-ordered products
\[
 \Theta_{\gamma_+}=\Sigma_2Z_{-3/4}\Sigma_1,
 \qquad
 \Theta_{\gamma_-}=\Sigma_1\Sigma_2
\]
are equal.
\item Both carry $W^{\min}$ to
\begin{equation}\label{eq:two-label-transported}
 W^{\mathrm{tr}}
 =x^4+y^4-2t_1x^5-4t_2x^4y+t_1t_2\bigl(8xy^5+4x^5y\bigr).
\end{equation}
\item The GKT transition \eqref{eq:mixed-GKT-transition} satisfies
\[
 \Sigma_2Z_{1/4}\Sigma_1=Z_1\Theta_{\gamma_+},
 \qquad
 Z_1W^{\mathrm{tr}}=W^{\max}.
\]
Thus the GKT transition and the path-ordered product of $\mathfrak D$ differ by $Z_1=\exp(t_1t_2X_{1,1})$.
\end{enumerate}
\end{proposition}

\begin{proof}
On $s_2=\varepsilon/10$ the three supports are met at
$s_1=-2\varepsilon/15$, $-\varepsilon/10$ and $-\varepsilon/15$, in the order
$\Sigma_1$-chord, mixed ray, $\Sigma_2$-chord.  The direction of travel $(1,0)$ has
positive pairing with $(3,4)$, $(1,1)$ and $(3,2)$, so
$\Theta_{\gamma_+}=\Sigma_2Z_{-3/4}\Sigma_1$.  On
$s_2=-\varepsilon/10$ the chords are met at $s_1=\varepsilon/15$ and
$2\varepsilon/15$, both positively, and the mixed ray is not met.  The
vertical segments lie on $s_1=\pm\varepsilon$ with $|s_2|\leq\varepsilon/10$,
where no support passes: the chords meet these lines at $|s_2|\geq
3\varepsilon/4$ and the ray at $s_2=\varepsilon$.  Hence
$\Theta_{\gamma_-}=\Sigma_1\Sigma_2$.

The formula in the proof of Theorem~\ref{thm:critical-embedding} gives
$[X_{1,0},X_{0,1}]=-3X_{1,1}$.  Modulo $(t_1^2,t_2^2)$ the commutator of the
logarithms of $\Sigma_2$ and $\Sigma_1$ is central, so
\[
 \Sigma_2\Sigma_1=\exp\!\left(\tfrac14t_1t_2[X_{0,1},X_{1,0}]\right)\Sigma_1\Sigma_2
 =Z_{3/4}\Sigma_1\Sigma_2.
\]
Therefore $\Theta_{\gamma_+}=Z_{-3/4}\Sigma_2\Sigma_1=\Sigma_1\Sigma_2=\Theta_{\gamma_-}$, which
proves~(1).

For (2), use \eqref{eq:two-label-derivatives} and $t_1^2=t_2^2=0$.  First,
\[
 \Sigma_1W^{\min}=x^4+y^4-2t_1x^5-2t_2y^5+2t_1t_2xy^5.
\]
The flow $Z_c$ changes only mixed terms, and
$Z_cV=V+8c\,t_1t_2(x^5y-xy^5)$ for every potential $V$ of the form
$x^4+y^4+O(t_1,t_2)$.  Hence
\[
 Z_{-3/4}\Sigma_1W^{\min}
 =x^4+y^4-2t_1x^5-2t_2y^5+t_1t_2\bigl(8xy^5-6x^5y\bigr).
\]
Applying $\Sigma_2$ adds $-4t_2x^4y+2t_2y^5$ from $x^4+y^4$ and $10t_1t_2x^5y$
from $-2t_1x^5$, which gives \eqref{eq:two-label-transported}.  Part~(1)
gives the same result along $\gamma_-$.

For (3), since $Z_c$ is central,
$\Sigma_2Z_{1/4}\Sigma_1=Z_{1/4+3/4}\Sigma_2Z_{-3/4}\Sigma_1=Z_1\Theta_{\gamma_+}$.  Finally
$Z_1W^{\mathrm{tr}}$ has mixed terms $t_1t_2\bigl((8-8)xy^5+(4+8)x^5y\bigr)$,
which is $W^{\max}$ by \eqref{eq:mixed-potentials}.
\end{proof}

\begin{remark}[Relation with the raywise factorization]
\label{rem:residual-not-correction}
Along $\gamma_+$ the mixed wall acts by its factor $Z_{-3/4}$, while the GKT
transition has $Z_{1/4}$ in the same position.
Proposition~\ref{prop:first-decoration} splits the wall factor as the GKT
factor times the residual $Z_{-1}$, and
Proposition~\ref{prop:two-label-transport}(3) is the corresponding statement
for the whole path: removing the residual turns the path-ordered product of
$\mathfrak D$ into the GKT transition.  By
Proposition~\ref{prop:two-label-transport}(2), $\mathfrak D$ itself does not
carry $W^{\min}$ to $W^{\max}$.
\end{remark}

\subsection{The maximal normal form}

For $a,b\in\QQ$ put
\[
 V(a,b)=x^4+y^4-2t_1x^5-4t_2x^4y+t_1t_2\bigl(a\,xy^5+b\,x^5y\bigr),
\]
so that $W^{\mathrm{tr}}=V(8,4)$ and $W^{\max}=V(0,12)$.  By the proof of
Proposition~\ref{prop:two-label-transport},
$Z_cV(a,b)=V(a-8c,b+8c)$.  Each orbit of the flows $Z_c$ in this family
therefore contains exactly one potential without an $xy^5$ term, and we
define
\begin{equation}\label{eq:endpoint-normal-form}
 \mathcal N^{\max}V(a,b)=Z_{a/8}V(a,b)=V(0,a+b).
\end{equation}
On the family $V(a,b)$ this operation is the maximal normal form of
Proposition~\ref{rem:normal-form-general} for the labels $i,j$; that normal
form keeps, for each support, only the balanced coefficient of greatest
$x$-degree.  The operation $\mathcal N^{\max}$ is idempotent and constant on
$Z$-orbits.
Deleting the $xy^5$ term instead would give $V(0,b)$, which differs from
$\mathcal N^{\max}V(a,b)$ whenever $a\neq0$.

\begin{proposition}[The normal form carries the incoming coefficient]
\label{prop:normal-form}
Let $W^{\min}_m$ be $W^{\min}$ with the coefficient $12$ of $t_1t_2xy^5$
replaced by $m\in\QQ$.  Then
\[
 \Theta_{\gamma_+}W^{\min}_m=V(m-4,4),
 \qquad
 \mathcal N^{\max}\Theta_{\gamma_+}W^{\min}_m=V(0,m).
\]
In particular $\mathcal N^{\max}\Theta_{\gamma_+}W^{\min}=W^{\max}$.
\end{proposition}

\begin{proof}
The operators $\Sigma_1$, $\Sigma_2$ and $Z_c$ fix the term $m\,t_1t_2xy^5$, and their
other contributions to the coefficient of $t_1t_2xy^5$ do not involve $m$.
The computation in the proof of
Proposition~\ref{prop:two-label-transport} with $12$ replaced by $m$
therefore gives $V(m-4,4)$, and \eqref{eq:endpoint-normal-form} gives
$V(0,m)$.
\end{proof}

The normal form transfers the mixed coefficient of $W^{\min}$ to $W^{\max}$;
it does not compute it.  The coefficient is fixed by periods,
which in turn do not select a representative in a $Z$-orbit.  Let
$W_0=x^4+y^4$ and $\Omega=dx\wedge dy$, and let $\nabla_0=\hbar d+dW_0\wedge$ act
on polynomial differential forms over $R_2[\hbar,\hbar^{-1}]$; this is the
operator of \eqref{eq:Brieskorn-module} over that ring.  Call a one-form \emph{relative} if its
pullbacks to both coordinate axes vanish, and call a two-form
\emph{relatively exact} if it equals $\nabla_0\beta$ for a relative one-form
$\beta$.

\begin{proposition}[Periods of the two-label family]
\label{prop:relative-exact}
\begin{enumerate}
\item The $t_1t_2$-part of $e^{(V(a,b)-W_0)/\hbar}\,\Omega$ is
$\nabla_0$-cohomologous to $\bigl(6-\tfrac{a+b}2\bigr)\,xy\,\Omega$.
\item The one-form $\beta=t_1t_2\,d(x^2y^2)$ is relative, and
$\nabla_0\beta=\bigl(W^{\max}-W^{\mathrm{tr}}\bigr)\Omega$; in particular
$\bigl(W^{\max}-W^{\mathrm{tr}}\bigr)\Omega$ is relatively exact.
\item Let $\hbar\in\CC^\times$ and let $\mathcal S\subset\CC^2$ be a properly
embedded oriented real surface whose boundary lies in the union of the
coordinate axes, on which $\operatorname{Re}(W_0/\hbar)\leq -c|z|^4+C$ for some $c>0$,
and whose area in the ball of radius $\rho$ grows at most polynomially in
$\rho$.  Then
\[
 \int_{\mathcal S} e^{W^{\max}/\hbar}\,\Omega
 =\int_{\mathcal S} e^{W^{\mathrm{tr}}/\hbar}\,\Omega
 \pmod{t_1^2,t_2^2}.
\]
\end{enumerate}
\end{proposition}

\begin{proof}
For a polynomial $f$, $\nabla_0(f\,dy)=(4x^3f+\hbar\,\partial_xf)\,\Omega$ and
$\nabla_0(f\,dx)=-(4y^3f+\hbar\,\partial_yf)\,\Omega$.  Taking $f=x^2y$, $xy^2$
and $x^6y$ gives the relations
$x^5y\,\Omega\sim-\tfrac\hbar2\,xy\,\Omega\sim xy^5\,\Omega$ and
$x^9y\,\Omega\sim\tfrac34\hbar^2\,xy\,\Omega$.  Modulo $(t_1^2,t_2^2)$ the
$t_1t_2$-part of $e^{(V(a,b)-W_0)/\hbar}$ is
$\hbar^{-1}(a\,xy^5+b\,x^5y)+\hbar^{-2}\,8x^9y$, which proves~(1).

The function $x^2y^2$ vanishes on both axes, so its differential pulls back
to zero there.  Since $d\beta=0$,
\[
 \nabla_0\beta=dW_0\wedge\beta
 =t_1t_2\bigl(4x^3dx+4y^3dy\bigr)\wedge\bigl(2xy^2dx+2x^2y\,dy\bigr)
 =t_1t_2\bigl(8x^5y-8xy^5\bigr)\Omega,
\]
which is $(W^{\max}-W^{\mathrm{tr}})\Omega$ by
\eqref{eq:two-label-transported} and \eqref{eq:mixed-potentials}.  Put
$\delta=W^{\max}-W^{\mathrm{tr}}\in t_1t_2\QQ[x,y]$.  Because
$W^{\mathrm{tr}}-W_0$ lies in $(t_1,t_2)$ and $t_1t_2(t_1,t_2)=0$,
\[
 e^{W^{\max}/\hbar}\Omega-e^{W^{\mathrm{tr}}/\hbar}\Omega
 =\hbar^{-1}e^{W_0/\hbar}\,\delta\,\Omega
 =\hbar^{-1}e^{W_0/\hbar}\nabla_0\beta
 =d\bigl(e^{W_0/\hbar}\beta\bigr).
\]
Under the stated growth conditions both integrals converge absolutely.  By
the coarea formula there is a sequence $\rho_n\to\infty$ of regular values of
$|z|$ on $\mathcal S$ and on $\partial\mathcal S$ for which the length of
$\mathcal S\cap\{|z|=\rho_n\}$ grows at most polynomially in $\rho_n$.  By
Stokes' theorem the integral of $d(e^{W_0/\hbar}\beta)$ over
$\mathcal S\cap\{|z|\le\rho_n\}$ equals the integral of
$e^{W_0/\hbar}\beta$ over $\partial\mathcal S\cap\{|z|\le\rho_n\}$, which
vanishes because $\beta$ is relative, plus an integral over
$\mathcal S\cap\{|z|=\rho_n\}$, which tends to zero because
$e^{W_0/\hbar}$ decays faster than the polynomial form $\beta$ and the length
grow.
\end{proof}

The class in Proposition~\ref{prop:relative-exact}(1) vanishes exactly when
$a+b=12$.  This is the GKT chamber condition for the support $\{i,j\}$,
which is a period condition \cite[Def.~3.47(3) and Thm.~4.9]{grossetal2022open}.
The periods of $V(a,b)$ depend only on $a+b$, and potentials in one $Z$-orbit
have the same periods.  Selecting the extremal representative within an
orbit therefore requires a normal-form rule; for these two labels, the rule
producing $W^{\max}$ is $\mathcal N^{\max}$.

\subsection{Two labels of the same twist}

The same phenomenon occurs without a mixed wall.  Take two labels with
parameters $t,t'$, both of twist $(1,0)$ and descendent degree one, and work
modulo $(t^2,t'^2)$.  Their singleton walls lie on the same chord and
commute, so modulo $(t^2,t'^2)$ the diagram $\mathfrak D$ has no further
wall, and every path crossing the chord in the direction of increasing slope
has path-ordered product
\[
 \Theta'=\exp\!\bigl(-\tfrac12(t+t')X_{1,0}\bigr).
\]
The support of the two labels has balanced profiles $(2,4)$ and $(6,0)$ and
critical exponent $(2,0)$.  Each label has the twist and descendent degree of
the label $i$, so the singleton terms are those of $t_1$ in
\eqref{eq:mixed-potentials}, and the extremal GKT potentials are
\[
 x^4+y^4-4(t+t')xy^4+p\,tt'x^2y^4
 \qquad\text{and}\qquad
 x^4+y^4-2(t+t')x^5+q\,tt'x^6
\]
for some $p,q\in\QQ$.  As in Proposition~\ref{prop:relative-exact}(1), the
GKT chamber condition for the support of the two labels is the vanishing of
the class of the $tt'$-part of $e^{(W-W_0)/\hbar}\,\Omega$ for each extremal
potential $W$.  Write $\sim$ for
$\nabla_0$-cohomology.  The formulas for $\nabla_0(f\,dx)$ with $f=x^2y$,
$x^2y^5$ and for $\nabla_0(f\,dy)$ with $f=x^3$, $x^7$ in the proof of
Proposition~\ref{prop:relative-exact} give
\[
 x^2y^4\,\Omega\sim-\tfrac{\hbar}4\,x^2\,\Omega,\qquad
 x^2y^8\,\Omega\sim\tfrac{5\hbar^2}{16}\,x^2\,\Omega,\qquad
 x^6\,\Omega\sim-\tfrac{3\hbar}4\,x^2\,\Omega,\qquad
 x^{10}\,\Omega\sim\tfrac{21\hbar^2}{16}\,x^2\,\Omega.
\]
Since $(t+t')^2=2tt'$ modulo $(t^2,t'^2)$, the $tt'$-parts for the two
potentials are $(\hbar^{-1}p\,x^2y^4+16\hbar^{-2}x^2y^8)\,\Omega$ and
$(\hbar^{-1}q\,x^6+4\hbar^{-2}x^{10})\,\Omega$, which are $\nabla_0$-cohomologous
to $(5-\frac p4)\,x^2\,\Omega$ and $(\frac{21}4-\frac{3q}4)\,x^2\,\Omega$.  The
chamber conditions therefore give $p=20$ and $q=7$, and the extremal
potentials are
\[
 W'^{\min}=x^4+y^4-4(t+t')xy^4+20\,tt'x^2y^4,
 \qquad
 W'^{\max}=x^4+y^4-2(t+t')x^5+7\,tt'x^6.
\]

\begin{proposition}[Same twist]
\label{prop:same-twist}
The diagram $\mathfrak D$ carries $W'^{\min}$ to
\[
 \Theta'W'^{\min}=x^4+y^4-2(t+t')x^5+tt'\bigl(6x^2y^4+5x^6\bigr),
\]
and the GKT transition is
$\Theta'\exp\!\bigl(\tfrac12tt'X_{2,0}\bigr)$.  Its extra factor lies on
the ray of shifted exponent $(3,1)$, which lies outside the cone spanned by
$(2,1)$ and $(1,2)$ and carries no wall of the diagram $\mathfrak D$.  The flows
$\exp(c\,tt'X_{2,0})$ change the coefficients $(p,q)$ of $tt'x^2y^4$ and
$tt'x^6$ to $(p-12c,q+4c)$, and the normal form without an $x^2y^4$ term
carries $\Theta'W'^{\min}$ to $W'^{\max}$.  If the coefficient $20$ is
replaced by $n$, the transported potential has coefficients $(n-14,5)$ and
its normal form has coefficient $(n+1)/3$ on $tt'x^6$.
\end{proposition}

\begin{proof}
Modulo $(t^2,t'^2)$ we have $(t+t')^2=2tt'$ and
$\Theta'=1-\tfrac12(t+t')X_{1,0}+\tfrac14tt'X_{1,0}^2$.  With
$X_{1,0}(xy^4)=-7x^2y^4$ and $X_{1,0}(x^5)=5x^6$, the three terms of
$\Theta'$ applied to $W'^{\min}$ give the displayed potential: the
coefficient of $tt'x^2y^4$ is $20-28+14=6$ and that of $tt'x^6$ is $5$;
with $n$ in place of $20$ the first becomes $n-14$.
Since $X_{2,0}(x^4+y^4)=4x^6-12x^2y^4$, the flow $\exp(c\,tt'X_{2,0})$
changes $(p,q)$ to $(p-12c,q+4c)$, and $c=\tfrac12$ carries $(6,5)$ to
$(0,7)$; in general $c=p/12$ gives $(0,q+p/3)$, which is $(0,(n+1)/3)$ for
$(p,q)=(n-14,5)$.  The flow of $tt'X_{2,0}$ commutes with $\Theta'$ modulo
$(t^2,t'^2)$.  The shifted exponent $(3,1)=c_1(2,1)+c_2(1,2)$ has
$c_2=-\tfrac13<0$.
\end{proof}

The path $\gamma_+$ meets the ray $L(3,1)$ before the $(2,1)$-chord.  In the
comparison of Section~\ref{sec:gkt-decoration} this ray has wall factor
$1$ and residual equal to the inverse of the GKT factor.

\subsection{Arbitrary labelled truncations}

\begin{proposition}[Maximal normal form]
\label{rem:normal-form-general}
Let $\Wcal$ be a labelled truncation (Definition~\ref{def:coefficient-window})
whose monomial set consists of square-free monomials, so that $R_{\Wcal,N}$ is
a quotient of $A_{\mathcal I}\otimes_\QQ\Bbbk$, with $A_{\mathcal I}$ the
labelled ring \eqref{eq:labelled-ring}, under $u_{i,1}\mapsto t_i$ and
$u_{i,d}\mapsto0$ for $d\neq1$.  For a support $J$, the \emph{diagonal} of $J$
consists of the coefficients of the monomials $t_Jx^{k_1}y^{k_2}$ with
$(k_1,k_2)$ one of the balanced profiles $b_{J,h}$ of
\eqref{eq:balanced-profiles}.  Let $W^{\min}$ and $W^{\max}$ denote the images
over $R_{\Wcal,N}$ of the extremal GKT potentials of the labels of $\Wcal$.
\begin{enumerate}
\item Let $V=x^4+y^4+\cdots$ be a potential over $R_{\Wcal,N}$ supported on
balanced monomials.  The orbit of $V$ under the group of the GKT critical
algebra of $\Wcal$ contains exactly one potential whose coefficients on the
diagonal of each support $J$ vanish except at $b_J^{\max}$.  We call it the
\emph{maximal normal form} $\mathcal N^{\max}V$ of $V$.
\item On the family $V(a,b)$, $\mathcal N^{\max}$ is the operation
\eqref{eq:endpoint-normal-form}; for the two labels of twist $(1,0)$ of
Proposition~\ref{prop:same-twist} it is the normal form used there.
\item The path-ordered product of $\mathfrak D$ along a path in the disk of
radius $2\varepsilon$ about $O$, such as $\gamma_+$, carries $W^{\min}$ to the
potential of a chamber index, and $\mathcal N^{\max}$ carries this potential
to $W^{\max}$.
\end{enumerate}
\end{proposition}

\begin{proof}
(1) The critical generators act triangularly on each diagonal
\cite[Thm.~1.3]{wuebben2026boundary}.  Hence the normal form exists and is
unique; it is constructed support by support in increasing $|J|$.

(2) The balanced profiles $b_J^{\max}$ of the supports $\{i\}$, $\{j\}$ and
$\{i,j\}$ are $(5,0)$, $(4,1)$ and $(5,1)$.  The potential
$V(0,a+b)=Z_{a/8}V(a,b)$ lies in the orbit of $V(a,b)$, and its diagonal
coefficients vanish except at these profiles, so it is
$\mathcal N^{\max}V(a,b)$ by (1).  For the two labels of
Proposition~\ref{prop:same-twist} the profiles $b_J^{\max}$ of the singletons
and of the pair are $(5,0)$ and $(6,0)$.  The normal form used there is
obtained by a flow $\exp(c\,tt'X_{2,0})$ of a critical generator, keeps the
singleton terms $-2(t+t')x^5$ and has no $tt'x^2y^4$ term, so it is
$\mathcal N^{\max}$ as well.

(3) Each wall factor of $\mathfrak D$ is the exponential of an element of the
critical algebra of $\Wcal$, so the path-ordered product along such a path
lies in its group and carries $W^{\min}$ to the potential of a chamber index
\cite[Thm.~4.27]{grossetal2022open}.  The maximal normal form of this
potential lies in the same orbit, so it is again the potential of a chamber
index, and its only nonzero balanced coefficients of nonempty support are
those at the profiles $b_J^{\max}$.  The only such chamber index is
$\nu^{\max}$ \cite[Thm.~1.3]{wuebben2026boundary},
\cite[Cor.~5.14]{grossetal2022open}, so $\mathcal N^{\max}$ carries the
transported $W^{\min}$ to $W^{\max}$.
\end{proof}

Propositions~\ref{prop:two-label-transport}, \ref{prop:normal-form}
and~\ref{prop:same-twist} make this correction explicit for two labels of
twists $(1,0)$ and $(0,1)$ and for two labels of twist $(1,0)$.

\begin{problem}[Maximal normalization]\label{prob:endpoint}
Let $\gamma$ be a path from $\ell$ to a region of $\mathfrak D$ adjacent to a
boundary circle of the annulus.  Construct, from the annular structure of
Theorem~\ref{thm:intro-reconstruction}, an operation $\mathcal E_\gamma$ on
potentials in the chart at the end of $\gamma$ such that, with
$\mathcal E_\gamma$ read in the frame at $\ell$ by transport along $\gamma$,
the operation $\Theta_\gamma^{-1}\mathcal E_\gamma\Theta_\gamma$ agrees with
$\mathcal N^{\max}$ on every truncation of
Proposition~\ref{rem:normal-form-general}.  For an intrinsic enumerative
construction, the construction should moreover single out the coefficients
of $W^{\min}$, for instance the value $m=12$ in
Proposition~\ref{prop:normal-form}, without using the GKT chamber conditions
of the supports involved.
\end{problem}

The comparisons above are local and
are made in one frame at $O$: the regions containing $\ell$ and $r$ are
adjacent to the joint, not to the boundary circles, and the computations see
neither the core holonomy $K$ nor the transitions at the singular points, so
a solution must be global.  A comparison theorem could also take the GKT
canonical systems as input, provided it identifies the boundary evaluation
with the normal form by a geometric argument.  Theta functions defined by
broken lines on the annulus, and a logarithmic compactification of the
finite stages (Section~\ref{sec:compactification}), are further open
problems.  The comparison between annular scattering and GKT open FJRW
theory proved here, Theorem~\ref{thm:intro-comparison} and
Proposition~\ref{thm:primary-Frobenius}, concerns wall factors, finite
homotopies and the primary Frobenius structure.
